problem:  
Let \((M,g)\) be a closed smooth \(n\)-dimensional Riemannian manifold.  
Fix \(\lambda>0\) and define the bistable potential \(V(u)=\tfrac14(u^2-1)^2\).  
Consider the stationary **Einstein–Allen–Cahn system**
\[
\begin{cases}
\operatorname{Ric}(g)=\lambda\big(\nabla u\otimes\nabla u+\tfrac14(u^2-1)^2 g\big),\\[0.3em]
\Delta_g u-(u^3-u)-\tfrac12R_g\,u+\alpha u=0,\\[0.3em]
\int_Mu^2\,d\mu_g=1,
\end{cases}
\]
where \(\alpha\in\mathbb{R}\) enforces the normalization constraint.

**Open problem (Einstein–Allen–Cahn rigidity conjecture):**  
Is every smooth solution \((g,u)\) of the above system necessarily *trivial*—that is, \(u\equiv\pm1\) and \(g\) Einstein—or can genuinely nonconstant “Allen–Cahn–Einstein” equilibria exist on closed manifolds for some dimension \(n\) and coupling \(\lambda>0\)?  

Equivalently, can one find (or rule out) a metric \(g\) and scalar field \(u\not\equiv\text{const.}\) on a compact manifold that jointly solve the coupled curvature–reaction system consistent with the Bianchi identity?

Why is it a “good” problem:  
This question is concise yet deep. It asks whether **coupled curvature–reaction equilibria** exist on compact manifolds for a fixed, explicit potential—analogous to an Einstein–scalar field system with Allen–Cahn self-interaction.  
No classical structure theorem resolves this case: the derivative coupling removes trivial proportionality, and any nonconstant solution would represent a new kind of **self-consistent geometric phase pattern**.  A proof of rigidity (all \(u\) constant) would extend Einstein’s uniqueness paradigm under nonlinear matter coupling, whereas a counterexample would reveal qualitatively new curvature–field configurations.  
Either outcome would connect geometric analysis, nonlinear elliptic PDE, and variational field theory, exemplifying a “narrow entrance, profound depth” problem: easy to state, formidable to resolve, and capable of bridging distinct areas of modern mathematics.